\documentclass[11pt,a4paper]{article}

\usepackage[T1]{fontenc}
\usepackage[utf8]{inputenc}
\usepackage{lmodern}
\usepackage{amssymb,amsmath,amsthm,mathtools}
\usepackage[margin=2.5cm]{geometry}
\usepackage{parskip}
\usepackage[hidelinks]{hyperref}
\usepackage{microtype}

% Theorem environments (arXiv / journal style)
\theoremstyle{plain}
\newtheorem{theorem}{Theorem}[section]
\newtheorem{proposition}[theorem]{Proposition}
\newtheorem{lemma}[theorem]{Lemma}
\newtheorem{corollary}[theorem]{Corollary}

\theoremstyle{definition}
\newtheorem{definition}[theorem]{Definition}
\newtheorem{example}[theorem]{Example}
\newtheorem{remark}[theorem]{Remark}

% Notation
\newcommand{\Ext}{\mathrm{Ext}}
\newcommand{\Shape}{\mathrm{Shape}}
\newcommand{\Pos}{\mathrm{Pos}}
\newcommand{\Type}{\mathbf{Type}}
\newcommand{\id}{\mathrm{id}}
\newcommand{\rootp}{\mathrm{o}}

\title{CoKleisli Structure and Strength\\
  for Directed Container Comonads}
\author{Extracted from a machine-checked development}
\date{}

\begin{document}
\maketitle

\begin{abstract}
We equip the comonad arising from a directed container with its coKleisli
operations and its canonical strength. Given a directed container, the
extension functor carries a comonad structure; we define coKleisli extension
and composition, prove the unit and associativity laws of the coKleisli
category, and construct the strength
$\sigma : W(A \times B) \to A \times W(B)$ together with its naturality and
counit laws. The proofs of all stated results have been machine-checked; we
record only the mathematical content.
\end{abstract}

\tableofcontents
\bigskip

%======================================================================
\section{Introduction}
%======================================================================

A directed container, in the sense of Ahman, Chapman, and Uustalu, is a
container equipped with a root choice, a sub-shape operation, and a shift on
positions, subject to five equational axioms (D1--D5). These data induce a
comonad on the extension of the underlying container. The coKleisli category
of that comonad is the natural home for ``contextual'' morphisms
$W(X) \to Y$, and the coKleisli extension operation is the unique lift of such
a morphism to an endomorphism of~$W$.

This note isolates the coKleisli layer and the canonical strength of the
comonad. We work throughout with a fixed directed container~$D$, writing $W$
for the extension of its underlying container, $\varepsilon$ for the counit
(extract at the root), and $\delta$ for the comultiplication (duplicate via
sub-shapes and shifts). No special categorical machinery beyond the notions of
comonad, coKleisli category, and strength is required. Proofs are indicated by
the phrase ``by lean,'' meaning they have been checked by a proof assistant
rather than written out in prose.

\paragraph{What this note claims.}
In ordinary mathematical language: for any directed container, coKleisli
extension satisfies the three standard laws
\[
  \mathrm{extend}\,\varepsilon = \id,\qquad
  \varepsilon \circ \mathrm{extend}\,f = f,\qquad
  \mathrm{extend}\,(g \circ \mathrm{extend}\,f)
    = \mathrm{extend}\,g \circ \mathrm{extend}\,f,
\]
and coKleisli composition is therefore a unital associative operation on
maps of the form $W(X) \to Y$. Moreover the strength
\[
  \sigma\langle s,\, v\rangle
  \;=\;
  \bigl( \pi_1\bigl(v(\rootp s)\bigr),\;
         \langle s,\, \pi_2 \circ v\rangle \bigr)
\]
is natural in both variables and satisfies $(\id \times \varepsilon) \circ \sigma
= \varepsilon$. All of this is formal content of the directed-container
comonad; nothing further is required.

%======================================================================
\section{Background: directed containers and their comonads}
%======================================================================

We recall only the data used below. A \emph{container} is a pair $(S,P)$
consisting of a type of shapes and a family of position types. Its
\emph{extension} at a type~$X$ is
\[
  \Ext\,C\,X \;=\; \sum_{s \,:\, S}\; \bigl(P(s) \to X\bigr).
\]
A \emph{directed container} equips a container with operations
\begin{align*}
  \rootp &\;:\; (s : S) \to P(s), \\
  \downarrow &\;:\; (s : S) \to P(s) \to S, \\
  \oplus &\;:\; (s : S) \to (p : P(s)) \to P(s \downarrow p) \to P(s),
\end{align*}
subject to five axioms D1--D5 (root-sub, root-shift, sub-root, nesting of
subs, and associativity of shift). These induce a comonad on $\Ext\,C$:
\begin{align*}
  \varepsilon\langle s,\, v\rangle &\;=\; v\bigl(\rootp s\bigr), \\
  \delta\langle s,\, v\rangle
  &\;=\;
  \bigl\langle s,\;
    \lambda p.\,
      \bigl\langle s \downarrow p,\;
        \lambda q.\, v(p \oplus q)\bigr\rangle
  \bigr\rangle.
\end{align*}
We write $W$ for this comonad (i.e.\ $W(X) = \Ext\,C\,X$) when the directed
container is fixed.

%======================================================================
\section{CoKleisli extension}
%======================================================================

\begin{definition}[CoKleisli extension]\label{def:extend}
Let $D$ be a directed container and let $f : W(X) \to Y$ be a function. The
\textbf{coKleisli extension} of~$f$ is the map
\begin{equation}
  \mathrm{extend}\,f \;:\; W(X) \;\longrightarrow\; W(Y)
\end{equation}
defined by
\begin{equation}
  \mathrm{extend}\,f\,\bigl\langle s,\, v \bigr\rangle
  \;=\;
  \bigl\langle s,\;
    \lambda p.\,
      f\bigl\langle
        s \downarrow p,\;
        \lambda q.\, v(p \oplus q)
      \bigr\rangle
  \bigr\rangle.
\end{equation}
\end{definition}

\begin{remark}
Definitionally, $\mathrm{extend}\,f = W(f) \circ \delta$: one first duplicates
the context and then applies~$f$ at every position. This is the standard
coKleisli extension of a comonad.
\end{remark}

\begin{proposition}[Extension of the counit]\label{prop:extend-counit}
For every type~$X$,
\begin{equation}
  \mathrm{extend}\,\varepsilon \;=\; \id_{W(X)}.
\end{equation}
\end{proposition}

\begin{proof}
By lean.
\end{proof}

\begin{remark}
Evaluating at $\langle s, v\rangle$ yields
$\langle s,\, \lambda p.\, v\bigl(p \oplus \rootp(s \downarrow p)\bigr)\rangle$,
which equals $\langle s, v\rangle$ by D3.
\end{remark}

\begin{proposition}[Counit after extension]\label{prop:counit-extend}
For every $f : W(X) \to Y$,
\begin{equation}
  \varepsilon \circ \mathrm{extend}\,f \;=\; f.
\end{equation}
\end{proposition}

\begin{proof}
By lean.
\end{proof}

\begin{remark}
This is the left unit law of the coKleisli category, and it characterises
$\mathrm{extend}\,f$ as the unique map $h$ with $\varepsilon \circ h = f$
compatible with the comultiplication in the usual way. The proof uses D1 and
D2: extracting after extension probes at the root sub-shape, which is the
original shape.
\end{remark}

\begin{proposition}[Extension respects composition]\label{prop:extend-extend}
For $f : W(X) \to Y$ and $g : W(Y) \to Z$,
\begin{equation}
  \mathrm{extend}\,(g \circ \mathrm{extend}\,f)
  \;=\;
  \mathrm{extend}\,g \circ \mathrm{extend}\,f.
\end{equation}
\end{proposition}

\begin{proof}
By lean.
\end{proof}

\begin{remark}
Both sides send $\langle s, v\rangle$ to a tree whose value at position~$p$ is
$g$ applied to a context built by nesting sub-shapes and shifts. Equality of
shapes uses D4; equality of values uses D5. This is the associativity law of
the coKleisli category, derived from coassociativity of~$\delta$.
\end{remark}

%======================================================================
\section{CoKleisli composition}
%======================================================================

\begin{definition}[CoKleisli composition]\label{def:cokleisli-comp}
Given coKleisli morphisms $f : W(X) \to Y$ and $g : W(Y) \to Z$, their
\textbf{coKleisli composite} is
\begin{equation}
  g \mathbin{\cdot_c} f \;:=\; g \circ \mathrm{extend}\,f
  \;:\; W(X) \;\longrightarrow\; Z.
\end{equation}
\end{definition}

\begin{proposition}[Right unit]\label{prop:cokleisli-right-unit}
For every $f : W(X) \to Y$,
\begin{equation}
  f \mathbin{\cdot_c} \varepsilon \;=\; f.
\end{equation}
\end{proposition}

\begin{proof}
By lean.
\end{proof}

\begin{proposition}[Left unit]\label{prop:cokleisli-left-unit}
For every $f : W(X) \to Y$,
\begin{equation}
  \varepsilon \mathbin{\cdot_c} f \;=\; f.
\end{equation}
\end{proposition}

\begin{proof}
By lean.
\end{proof}

\begin{proposition}[Associativity]\label{prop:cokleisli-assoc}
For $f : W(X) \to Y$, $g : W(Y) \to Z$, and $h : W(Z) \to U$,
\begin{equation}
  \bigl(h \mathbin{\cdot_c} g\bigr) \mathbin{\cdot_c} f
  \;=\;
  h \mathbin{\cdot_c} \bigl(g \mathbin{\cdot_c} f\bigr).
\end{equation}
\end{proposition}

\begin{proof}
By lean.
\end{proof}

Taken together, Propositions~\ref{prop:cokleisli-right-unit}--\ref{prop:cokleisli-assoc}
say that coKleisli morphisms form a category with identities given by the
components of~$\varepsilon$.

%======================================================================
\section{Canonical strength}
%======================================================================

Every comonad on a cartesian category is canonically strong. For a directed
container comonad the strength admits a direct combinatorial description.

\begin{definition}[Strength]\label{def:strength}
The \textbf{canonical strength} of the directed container comonad is the
family of maps
\begin{equation}
  \sigma_{A,B} \;:\; W(A \times B) \;\longrightarrow\; A \times W(B)
\end{equation}
defined by
\begin{equation}
  \sigma_{A,B}\,\bigl\langle s,\, v \bigr\rangle
  \;=\;
  \Bigl(
    \pi_1\bigl(v(\rootp s)\bigr),\;
    \bigl\langle s,\; \pi_2 \circ v \bigr\rangle
  \Bigr).
\end{equation}
The first component extracts the $A$-part at the root; the second retains the
shape and projects every label to its $B$-part.
\end{definition}

\begin{proposition}[Naturality in the left variable]\label{prop:strength-nat-left}
For every $f : A \to A'$,
\begin{equation}
  (f \times \id) \circ \sigma_{A,B}
  \;=\;
  \sigma_{A',B} \circ W(f \times \id).
\end{equation}
\end{proposition}

\begin{proof}
By lean.
\end{proof}

\begin{proposition}[Naturality in the right variable]\label{prop:strength-nat-right}
For every $g : B \to B'$,
\begin{equation}
  (\id \times W(g)) \circ \sigma_{A,B}
  \;=\;
  \sigma_{A,B'} \circ W(\id \times g).
\end{equation}
\end{proposition}

\begin{proof}
By lean.
\end{proof}

\begin{proposition}[Strength and counit]\label{prop:strength-counit}
\begin{equation}
  (\id \times \varepsilon) \circ \sigma_{A,B}
  \;=\;
  \varepsilon_{A \times B}.
\end{equation}
\end{proposition}

\begin{proof}
By lean.
\end{proof}

\begin{remark}
Extracting the $B$-component of the strength is the same as extracting
directly from the $A \times B$-labelled tree: both pick the $B$-part of the
label at the root.
\end{remark}

%======================================================================
\section{Discussion}
%======================================================================

\paragraph{What has been established.}
The coKleisli category of a directed container comonad has been constructed
explicitly from root, sub, and shift. The unit laws rest on D1--D3; the
associativity law rests on D4--D5. The strength is the standard one for a
comonad on a cartesian category, written out in container coordinates.

\paragraph{Scope.}
The development treats a fixed directed container and does not re-derive the
comonad laws themselves (those belong to the underlying directed-container
theory). Strength is stated only with the counit interaction recorded above;
further strength axioms (e.g.\ interaction with comultiplication) are not
listed separately here.

\paragraph{Machine checking.}
Every definition and theorem statement above corresponds to a declaration in a
machine-checked formalisation. The phrase ``by lean'' indicates that the proof
has been verified by a proof assistant; we have not reproduced the formal
proof scripts here.

%======================================================================
\section*{Acknowledgements}
%======================================================================

The mathematical content follows the theory of directed containers of Ahman,
Chapman, and Uustalu, and standard coKleisli calculus for comonads. The
presentation here is a prose rendering of a formal development intended for
human review of the theorem statements.

\begin{thebibliography}{9}

\bibitem{ahman2014}
D.~Ahman, J.~Chapman, and T.~Uustalu.
\newblock When is a container a comonad?
\newblock \emph{Logical Methods in Computer Science}, 10(3), 2014.

\bibitem{abbott2005}
M.~Abbott, T.~Altenkirch, and N.~Ghani.
\newblock Containers: constructing strictly positive types.
\newblock \emph{Theoretical Computer Science}, 342(1):3--27, 2005.

\end{thebibliography}

\end{document}
